\documentclass[11pt, parskip=half]{scrbook}

\usepackage{azuma}
\usepackage{stablespacing}
\usepackage{minted}
\usepackage{codeblock}
\usepackage[a4paper, margin=1in]{geometry} % center margins

\usepackage{enumitem}
\setlist[itemize]{topsep=0pt, itemsep=0.1em} %partopsep, parsep

\raggedbottom

\begin{document}
\title{Some title}
\maketitle
\tableofcontents

% part, chapter, section, subsection, subsubsection, paragraph, subparagraph
\part{some part}
\chapter{some chapter}
\section{some section}
\subsection{some subsection}
\subsubsection{some subsubsection}
\paragraph{some paragraph}
\subparagraph{some subparagraph}


% minted example
\begin{minted}{python}
    def incmatrix(genl1,genl2):
    m = len(genl1)
    n = len(genl2)
    M = None #to become the incidence matrix
    VT = np.zeros((n*m,1), int)  #dummy variable
    
    #compute the bitwise xor matrix
    M1 = bitxormatrix(genl1)
    M2 = np.triu(bitxormatrix(genl2),1) 

    for i in range(m-1):
        for j in range(i+1, m):
            [r,c] = np.where(M2 == M1[i,j])
            for k in range(len(r)):
                VT[(i)*n + r[k]] = 1;
                VT[(i)*n + c[k]] = 1;
                VT[(j)*n + r[k]] = 1;
                VT[(j)*n + c[k]] = 1;
                
                if M is None:
                    M = np.copy(VT)
                else:
                    M = np.concatenate((M, VT), 1)
                
                VT = np.zeros((n*m,1), int)
    return M
\end{minted}

% same example but with my styling, different environment
\begin{codeblock}{python}
    def incmatrix(genl1,genl2):
    m = len(genl1)
    n = len(genl2)
    M = None #to become the incidence matrix
    VT = np.zeros((n*m,1), int)  #dummy variable
    
    #compute the bitwise xor matrix
    M1 = bitxormatrix(genl1)
    M2 = np.triu(bitxormatrix(genl2),1) 

    for i in range(m-1):
        for j in range(i+1, m):
            [r,c] = np.where(M2 == M1[i,j])
            for k in range(len(r)):
                VT[(i)*n + r[k]] = 1;
                VT[(i)*n + c[k]] = 1;
                VT[(j)*n + r[k]] = 1;
                VT[(j)*n + c[k]] = 1;
                
                if M is None:
                    M = np.copy(VT)
                else:
                    M = np.concatenate((M, VT), 1)
                
                VT = np.zeros((n*m,1), int)
    return M

\end{codeblock}

\end{document}
